\documentclass[10pt,a4paper]{amsart}
\usepackage[margin=19mm]{geometry}
\usepackage{amsmath,amssymb,mathtools}
\usepackage[scr=rsfs,cal=euler]{mathalfa}
\usepackage{xfrac}
\usepackage[unicode,hidelinks,bookmarks=false]{hyperref}
\newcommand{\ie}{i.e.\ }
\newcommand{\cf}{cf.\ }
\newcommand{\resp}{respectively\ }
\newcommand{\Subsection}{\subsection}
\providecommand{\nobreakdash}{\nobreak\hbox{-}}
\setlength{\emergencystretch}{2em}
\multlinegap0pt
\usepackage{bm}
\usepackage{bbm}
\usepackage{dsfont}
\usepackage{xspace}
\usepackage{cancel}
\usepackage{mathtools}
\usepackage[shortlabels]{enumitem}
\usepackage{amsthm}
\makeatletter
\@ifundefined{theorem}{\newtheorem{theorem}{Theorem}}{}
\@ifundefined{definition}{\newtheorem{definition}[theorem]{Definition}}{}
\makeatother
\title[Local Antisymmetric Connectedness in Quasi-Uniform and Quasi]{Local Antisymmetric Connectedness in Quasi-Uniform and Quasi-Modular Spaces}
\author{Philani Rodney Majozi}
\date{Source: arXiv:2601.16361v1 (CC BY 4.0)}
\begin{document}
\maketitle

\noindent\emph{Source and licence.} Local Antisymmetric Connectedness in Quasi-Uniform and Quasi-Modular Spaces. Version arXiv:2601.16361v1 (CC BY 4.0), distributed under the Creative Commons Attribution 4.0 International licence, as recorded in the arXiv licence metadata for this exact version and shown on the abstract page https://arxiv.org/abs/2601.16361v1. Full rights notice: licenses/arxiv-2601.16361-CC-BY-4.0.txt.\par\smallskip

\noindent\textit{Editorial scope.} Counted unit: the Theorem whose source label is \texttt{thm:quni-characterization} in the article (source0.tex lines 1116--1124, source label \texttt{thm:quni-characterization}), delivered together with the article's own complete original proof of it (source0.tex lines 1126--1169, one \texttt{proof} environment of 44 lines). No mathematics is written, repaired or re-derived: statement and proof are the article's own text line for line. Required context retained uncounted: def:qu:basic (lines 311--346). Every definition, display and label of those lines is retained unchanged. Omitted from this package and not redistributed: 1--310, 347--1115, 1125--1125, 1170--1951. Nothing inside a retained excerpt is omitted or altered: every retained body is delivered exactly as the article's lines are written, so no omission receipt of any kind accompanies this package. Citations: the retained excerpts carry no citation command at all, so no bibliography entry of the article is transcribed and no reference is redistributed. Reference adaptation: none was needed and none was made; every reference inside the delivered unit resolves inside it, so no unresolved reference or citation is produced. Printed slips kept literal: no printed word, symbol, hypothesis or number of the counted statement or of its proof was altered. Layout and class: the article's own class is \texttt{elsarticle} with packages including amsmath, amssymb, amsthm, enumitem, geometry, hyperref, inputenc, mathtools; the wrapper is the lane's pinned \texttt{amsart} editorial wrapper at 10pt on A4, which restates the theorem-like environments the retained text needs and adds the article's own macro definitions that the retained text needs (none), with extra wrapper packages (bm, bbm, dsfont, xspace, cancel, mathtools, enumitem). The delivered numbering is the unit's own. Assets: no retained span contains a figure environment, a graphics-inclusion command, a PDF-page inclusion, a table or a TikZ picture; no image file is copied into this package, the package declares no asset dependency and no image file is redistributed with it. Licence: version arXiv:2601.16361v1 of this article is distributed under the Creative Commons Attribution 4.0 International licence, as recorded in the arXiv licence metadata for this exact version and shown on the abstract page https://arxiv.org/abs/2601.16361v1. A full rights notice is delivered at licenses/arxiv-2601.16361-CC-BY-4.0.txt. Characters: every retained excerpt is pure ASCII, so the delivered engine, XeLaTeX with its default Latin Modern text font, reports no missing character.
\begin{definition}
\label{def:qu:basic}
Let $X$ be a set and let
$\Delta=\{(x,x)\mid x\in X\}\subseteq X\times X$. A \emph{quasi-uniformity}
on $X$ is a filter $\mathcal{U}$ on $X\times X$ such that
\begin{enumerate}[label=\textup{(\roman*)}]
\item $\Delta\subseteq U$ for every $U\in\mathcal{U}$,
\item for every $U\in\mathcal{U}$ there exists $V\in\mathcal{U}$ such
that $V\circ V\subseteq U$, where
$R\circ S=\{(x,z):\exists y,\ (x,y)\in R,\ (y,z)\in S\}$.
\end{enumerate}
The pair $(X,\mathcal{U})$ is called a \emph{quasi-uniform space}, and
the members of $\mathcal{U}$ are called \emph{entourages}. The
\emph{conjugate} quasi-uniformity is
\[
\mathcal{U}^{-1}:=\{U^{-1}\mid U\in\mathcal{U}\},
\qquad
U^{-1}:=\{(y,x):(x,y)\in U\}.
\]
The \emph{associated supremum uniformity} is the coarsest uniformity
finer than both $\mathcal{U}$ and $\mathcal{U}^{-1}$, denoted
\[
\mathcal{U}^{*}:=\mathcal{U}\vee\mathcal{U}^{-1}.
\]
For $U\in\mathcal{U}$ and $x\in X$ set $U(x):=\{y\in X:(x,y)\in U\}$.
The topology induced by $\mathcal{U}$ is
\[
\tau(\mathcal{U})
:=\{O\subseteq X:\forall x\in O\ \exists U\in\mathcal{U}\ \text{with }
U(x)\subseteq O\}.
\]
Analogously $\tau(\mathcal{U}^{-1})$ is induced by $\mathcal{U}^{-1}$,
and $(X,\tau(\mathcal{U}),\tau(\mathcal{U}^{-1}))$ is the induced
\emph{bitopological space}. We also write $\tau(\mathcal{U}^{*})$ for
the uniform topology of $\mathcal{U}^{*}$.
\end{definition}

\begin{theorem}
\label{thm:quni-characterization}
Let $(X,\mathcal U^{+})$ be a quasi-uniform space and write
$\tau^{+}=\tau(\mathcal U^{+})$, $\tau^{-}=\tau((\mathcal U^{+})^{-1})$,
and $\tau^\vee=\tau^{+}\vee\tau^{-}$ (Definition~\ref{def:qu:basic}).
Then $X$ is locally antisymmetrically connected if and only if every
point admits a base of entourages whose induced $\tau^\vee$-neighborhoods
are antisymmetrically connected.
\end{theorem}

\begin{proof}
Assume first that $X$ is locally antisymmetrically connected.
Fix $x\in X$ and let $U$ be an arbitrary $\tau^\vee$-neighborhood of $x$.
By the definition of $\tau^\vee$, there exist entourages
$E,F\in\mathcal U^{+}$ such that
\[
E(x)\cap F^{-1}(x)\subseteq U.
\]
Since $X$ is locally antisymmetrically connected at $x$, there exists a
$\tau^\vee$-neighborhood $V$ of $x$ satisfying
\[
x\in V\subseteq E(x)\cap F^{-1}(x),
\]
with $V$ antisymmetrically connected.
Again by the definition of $\tau^\vee$, there exist entourages
$G,H\in\mathcal U^{+}$ such that
\[
G(x)\cap H^{-1}(x)\subseteq V.
\]
Then $G(x)\cap H^{-1}(x)$ is a $\tau^\vee$-neighborhood of $x$ which is
antisymmetrically connected. Since $U$ was arbitrary, such neighborhoods
form a base at $x$.

Conversely, assume that for each $x\in X$ there exists a base
$\mathcal B_x\subseteq\mathcal U^{+}$ such that for every
$E\in\mathcal B_x$ the set
\[
E(x)\cap E^{-1}(x)
\]
is antisymmetrically connected.
Let $U$ be any $\tau^\vee$-neighborhood of $x$.
Then there exist entourages $F,G\in\mathcal U^{+}$ with
\[
F(x)\cap G^{-1}(x)\subseteq U.
\]
Choose $E\in\mathcal B_x$ with $E\subseteq F\cap G$.
Then
\[
E(x)\cap E^{-1}(x)\subseteq F(x)\cap G^{-1}(x)\subseteq U,
\]
and by assumption $E(x)\cap E^{-1}(x)$ is antisymmetrically connected.
Hence $X$ is locally antisymmetrically connected at $x$.
Since $x$ was arbitrary, $X$ is locally antisymmetrically connected.
\end{proof}

\end{document}
